\documentclass[11pt]{article}
\usepackage[margin=1in]{geometry}
\usepackage[T1]{fontenc}
\usepackage[utf8]{inputenc}
\usepackage{booktabs}
\usepackage{hyperref}
\usepackage{array}
\usepackage{longtable}

\title{MCQ Contract Audit for the Qwen3-1.7B Rollout Bundle}
\date{2026-03-16 UTC}

\begin{document}
\maketitle

\section*{Bottom Line}
The current \texttt{GPQA} and \texttt{MMLU-Pro} numbers in the rollout report are
not trustworthy as multiple-choice accuracy estimates. A bounded live audit on the
same prompt/decode contract shows that the current scorer materially undercounts
answer-bearing outputs because it requires the final non-empty line to be exactly
\texttt{Answer: X}. The same raw completions frequently end with boxed letters,
\texttt{Final Answer:}, or freeform tail text like ``the answer is I'', all of which
the old scorer marks wrong.

The \texttt{LiveCodeBench} concern is separate: its recovered number is still useful as
task-local telemetry, but it is not benchmark-comparable because the current run uses
a question-id-sorted 800-problem slice plus a local prompt/style mapping rather than
the benchmark-native protocol.\footnote{\url{https://github.com/LiveCodeBench/LiveCodeBench}}

\section*{Subset Audit}
I ran two bounded one-GPU audits under the \emph{old} MC contract and saved the raw
completions to
\texttt{outputs/weeks/2026-W12/mcq\_contract\_audit\_2026-03-16/}. The audits used
\texttt{Qwen/Qwen3-1.7B}, \texttt{temperature=0.2}, \texttt{num\_generations=10},
\texttt{max\_tokens=4096}, and the same chat-template MC prompts as the shipped
bundle. Qwen's own model card recommends standardizing MC evaluation with a JSON
\texttt{answer} field rather than an ad hoc terminal string.\footnote{\url{https://huggingface.co/Qwen/Qwen3-1.7B}}

\begin{center}
\begin{tabular}{lrrrrrr}
\toprule
Dataset & Prompts & Rollouts & Strict Correct & Obvious Correct & Strict Stop Correct & Obvious Stop Correct \\
\midrule
GPQA subset & 16 & 160 & 11 & 24 & 11/39 & 20/39 \\
MMLU-Pro subset & 8 & 80 & 9 & 23 & 9/50 & 16/50 \\
\bottomrule
\end{tabular}
\end{center}

\noindent
``Strict'' means the old collector policy: only accept a terminal line exactly equal
to \texttt{Answer: X}. ``Obvious'' means a bounded tail heuristic over the same raw
outputs that recovers boxed answers, \texttt{Answer: X} variants, and ``the answer is
X'' endings. This is intentionally conservative; it still undercounts some long
truncated outputs that visibly commit to an option before running on.

\section*{Representative Misses}
\textbf{GPQA stop completions mis-scored as wrong under the old parser.}
\begin{itemize}
\item A stopped rollout ended with \texttt{\textbackslash boxed\{D\}} after a
\texttt{Final Answer:} header. Gold was \texttt{D}; strict score was wrong.
\item Another stopped rollout ended with
\texttt{\textbackslash boxed\{\textbackslash text\{Answer: A\}\}}. Gold was
\texttt{A}; strict score was wrong.
\item A length-truncated rollout explicitly said ``the answer is B'' before continuing
to revisit the options. Gold was \texttt{B}; strict score was wrong.
\end{itemize}

\textbf{MMLU-Pro shows the same pathology.}
\begin{itemize}
\item A stopped rollout ended with \texttt{\textbackslash boxed\{I\}} after an
\texttt{Answer:} header. Gold was \texttt{I}; strict score was wrong.
\item Another stopped rollout ended with
\texttt{\textbackslash boxed\{\textbackslash text\{F\}\}}. Gold was \texttt{F};
strict score was wrong.
\item Multiple length-truncated rollouts explicitly said ``the answer is I'' in the
tail while the strict parser still returned no answer.
\end{itemize}

\section*{Interpretation}
Two failure modes are active at once.
\begin{enumerate}
\item \textbf{Parser loss is real and material.} Even among \emph{stopped} outputs,
the old parser leaves many answer-bearing generations on the floor. On the bounded
GPQA subset it recovers only 11 obviously correct stopped generations, while the tail
audit finds 20. On the bounded MMLU subset it recovers only 9, while the tail audit
finds 16.
\item \textbf{The prompt/decode contract still produces long unfinished thought
traces.} Many subset rollouts also hit the 4096-token forensic cap before reaching a
clean final line, so the old parser is not the entire story. But the audit proves the
old MC report is not merely ``low because the model failed''; a nontrivial slice of
the collapse is scoring-contract loss.
\end{enumerate}

\section*{Code Repair}
The project code is patched locally but not yet rerun on the full bundle.
\begin{itemize}
\item GPQA and MMLU-Pro prompts now request a final JSON object
\texttt{\{"answer": "X"\}} rather than the brittle \texttt{Answer: X} line.
\item The MC parser now reads structured tail answers more robustly, including JSON
\texttt{answer} fields, \texttt{Final Answer:}, boxed letters, and bare terminal
choice letters in the final lines.
\item The report builder text has been updated so future rebuilt PDFs describe the new
MC contract rather than the old one.
\end{itemize}

\section*{Implications for the Existing Report}
\begin{itemize}
\item Treat the existing \texttt{GPQA} and \texttt{MMLU-Pro} accuracy claims as
invalid under the old MC contract.
\item Treat the recovered \texttt{LiveCodeBench} number as local telemetry, not a
benchmark-comparable claim.
\item The math datasets are unaffected by this specific MC parsing bug.
\item The next trustworthy step is to rerun \texttt{GPQA} and \texttt{MMLU-Pro} under
the patched MC contract, then decide separately whether \texttt{LiveCodeBench} should
also be rerun under a benchmark-native representative protocol.
\end{itemize}

\end{document}
